% Aqui começa o capítulo de Análise de Riscos e Incertezas.
\chapter{Análise de Riscos e Incertezas}\label{chp:RiscosIncertezas}

Projetos de tecnologia aplicados à área social estão inseridos em ambientes caracterizados por elevada complexidade, múltiplos interesses e mudanças frequentes de contexto. Nesse cenário, o gerenciamento de riscos e incertezas não pode ser tratado como uma atividade pontual ou meramente documental, mas como um processo contínuo de identificação, análise e resposta, em alinhamento com os princípios e domínios de desempenho do PMBOK 7ª edição.

Neste capítulo, são analisados os principais riscos e fontes de incerteza associados ao desenvolvimento do sistema de gestão de casos sociais, bem como as estratégias adotadas para seu tratamento no contexto de uma abordagem ágil baseada em Scrum.

\section{Natureza da Incerteza no Projeto}

A incerteza no projeto manifesta-se de diversas formas, incluindo requisitos incompletos ou mutáveis, comportamento imprevisível dos usuários, dependência de decisões institucionais e limitações de recursos. No caso específico de um sistema de gestão de casos sociais, a diversidade de situações atendidas e a sensibilidade das informações envolvidas ampliam esse grau de incerteza.

Diferentemente de abordagens tradicionais, que buscam eliminar a incerteza por meio de planejamento detalhado antecipado, a abordagem adotada neste projeto reconhece a incerteza como um elemento inerente e busca gerenciá-la por meio de ciclos curtos de aprendizado, validação contínua e adaptação progressiva.

\section{Identificação dos Principais Riscos}

A identificação de riscos foi realizada de forma iterativa, considerando tanto aspectos técnicos quanto organizacionais e humanos. Entre os principais riscos identificados destacam-se:

\begin{itemize}
    \item Baixa adoção do sistema pelos usuários finais, devido à resistência à mudança ou inadequação da solução ao fluxo de trabalho existente.
    \item Inconsistência ou baixa qualidade dos dados registrados, comprometendo a confiabilidade das análises e relatórios.
    \item Mudanças nos requisitos decorrentes de alterações institucionais ou normativas.
    \item Limitações de tempo e recursos humanos disponíveis para o desenvolvimento do projeto.
    \item Riscos relacionados à segurança da informação e à privacidade dos dados dos casos sociais.
\end{itemize}

Esses riscos foram priorizados com base em sua probabilidade de ocorrência e impacto potencial sobre os objetivos do projeto, considerando o contexto institucional simulado.

\section{Estratégias de Resposta aos Riscos}

As estratégias de resposta aos riscos adotadas no projeto estão alinhadas à lógica ágil de mitigação contínua, priorizando a redução da incerteza por meio da entrega incremental e do feedback frequente.

Para mitigar o risco de baixa adoção do sistema, são realizadas validações frequentes com usuários e stakeholders durante as revisões de sprint, permitindo ajustes precoces na usabilidade e no fluxo das funcionalidades. A qualidade dos dados é tratada por meio da definição de critérios claros de cadastro e da incorporação gradual de mecanismos de validação.

Mudanças de requisitos são absorvidas por meio do refinamento contínuo do Product Backlog, evitando comprometer o planejamento de curto prazo das sprints. As limitações de tempo e recursos são tratadas por meio da priorização rigorosa do backlog, concentrando esforços nas funcionalidades de maior valor.

No que se refere à segurança e privacidade, o projeto incorpora preocupações éticas desde as fases iniciais, influenciando decisões sobre controle de acesso, perfis de usuário e tratamento das informações sensíveis.

\section{Monitoramento Contínuo dos Riscos}

O monitoramento dos riscos ocorre de forma contínua ao longo do projeto, integrado às cerimônias do Scrum. Durante as reuniões diárias e as retrospectivas de sprint, a equipe avalia a emergência de novos riscos, a eficácia das respostas adotadas e a necessidade de ajustes no planejamento.

Essa abordagem permite tratar riscos como elementos dinâmicos, que evoluem ao longo do tempo, em vez de itens estáticos registrados em um documento formal. O desempenho nesse aspecto está relacionado à capacidade da equipe de antecipar problemas, aprender com os erros e adaptar o projeto de forma consciente.

\section{Riscos, Incerteza e Tomada de Decisão}

A análise de riscos e incertezas desempenha papel central na tomada de decisão ao longo do projeto. Decisões relacionadas à priorização do backlog, ao sequenciamento das entregas e à alocação de esforços são influenciadas pela avaliação contínua dos riscos mais relevantes.

Dessa forma, o gerenciamento de riscos não é tratado como uma atividade isolada, mas como um componente integrado ao planejamento ágil e à governança do projeto, em consonância com os domínios de desempenho e os princípios do PMBOK 7ª edição.
